%%
%% \BibTeX command to typeset BibTeX logo in the docs
\AtBeginDocument{%
  \providecommand\BibTeX{{%
    Bib\TeX}}}

%% Rights management information.  This information is sent to you
%% when you complete the rights form.  These commands have SAMPLE
%% values in them; it is your responsibility as an author to replace
%% the commands and values with those provided to you when you
%% complete the rights form.
\setcopyright{acmcopyright}
\copyrightyear{2018}
\acmYear{2018}
\acmDOI{XXXXXXX.XXXXXXX}

%% These commands are for a PROCEEDINGS abstract or paper.
\acmConference[Conference acronym 'XX]{Make sure to enter the correct
  conference title from your rights confirmation emai}{June 03--05,
  2018}{Woodstock, NY}
\acmISBN{978-1-4503-XXXX-X/18/06}

% %% The abstract is a short summary of the work to be presented in the
% %% article.
% %% The following content must be adapted for the final version
% % paper-specific
% \newcommand\vldbdoi{XX.XX/XXX.XX}
% \newcommand\vldbpages{XXX-XXX}
% % issue-specific
% \newcommand\vldbvolume{14}
% \newcommand\vldbissue{1}
% \newcommand\vldbyear{2020}
% % should be fine as it is
% \newcommand\vldbauthors{\authors}
% \newcommand\vldbtitle{\shorttitle} 
% % leave empty if no availability url should be set
% \newcommand\vldbavailabilityurl{URL_TO_YOUR_ARTIFACTS}
% % whether page numbers should be shown or not, use 'plain' for review versions, 'empty' for camera ready
% \newcommand\vldbpagestyle{plain} 


\usepackage{xspace}
\setlength{\emergencystretch}{2em}
\hbadness=10000
\vbadness=10000
\hfuzz=1pt
\vfuzz=3pt

%%%% self import %%%%%
\usepackage[normalem]{ulem}
% \useunder{\uline}{\ul}{}
\usepackage{amsmath,amsfonts}
    \usepackage{algorithmic}
\usepackage{multirow}
\usepackage{graphicx, caption, subcaption}
\usepackage{tabularx}
\usepackage{textcomp}

\usepackage{soul} % to allow \ul
\usepackage{amsfonts}
% \usepackage{subcaption}
%%% CODE %%% 
\usepackage{diagbox}
\usepackage[linesnumbered,ruled,vlined]{algorithm2e}
% \usepackage[dvipsnames]{xcolor}
% Definindo novas cores
\definecolor{verde}{rgb}{0.25,0.5,0.35}
\definecolor{jpurple}{rgb}{0.5,0,0.35}
\definecolor{darkgreen}{rgb}{0.0, 0.2, 0.13}
% Configurando layout para mostrar codigos Java
\usepackage{listings}
\newcommand{\estiloJava}{
\lstset{
    language=Java,
    basicstyle=\ttfamily\scriptsize,
    keywordstyle=\color{jpurple}\bfseries,
    stringstyle=\color{red},
    commentstyle=\color{verde},
    morecomment=[s][\color{blue}]{/**}{*/},
    extendedchars=true,
    showspaces=false,
    showstringspaces=false,
    numbers=left,
    numberstyle=\tiny,
    breaklines=true,
    backgroundcolor=\color{cyan!10},
    breakautoindent=true,
    captionpos=b,
    xleftmargin=0pt,
    tabsize=2,
}}
\lstset{
    showstringspaces=false,
    basicstyle=\ttfamily,
    keywordstyle=\color{blue},
    commentstyle=\color[blue]{0.9},
    stringstyle=\color[RGB]{255,150,75},
    morekeywords={Event, EventBlotter}
}

\definecolor{almond}{rgb}{0.98, 0.98, 0.82}
\definecolor{blue}{rgb}{0.0, 0.0, 1.0}
\newcommand{\inlinecode}[2]{\colorbox{almond}{\lstinline[language=#1]$#2$}}

\usepackage{tikz}
\newcommand*\circled[1]{\tikz[baseline=(char.base)]{
            \node[shape=circle,fill,inner sep=1pt] (char) {\textcolor{white}{#1}};}}

\usepackage{enumitem}
\newenvironment{myitemize}
{ \begin{itemize}[leftmargin=0.2in]	
		\vspace{-1ex}	
		\setlength{\itemsep}{0pt}
		\setlength{\parskip}{0pt}
		\setlength{\parsep}{0pt}    }
	{ 	 \end{itemize}                    }

\newenvironment{myenumerate}
{ \begin{enumerate}[leftmargin=0.2in]
		\vspace{-1ex}
		\setlength{\itemsep}{0pt}
		\setlength{\parskip}{0pt}
		\setlength{\parsep}{0pt}    }
	{ \end{enumerate}                  }
	
% \newcommand{\huilin}[1]{{\textcolor{blue}{#1}}}            
% \newcommand{\comment}[1]{{}}
\newcommand{\tony}[1]{{\textcolor{red}{#1}}}
\newcommand{\borrow}[1]{{\textcolor{purple}{#1}}}
\newcommand{\wangxin}[1]{{\textcolor{orange}{#1}}}
\newcommand{\modify}[1]{{\textcolor{violet}{#1}}}
\newcommand{\xianzhi}[1]{{\textcolor{yellow}{#1}}}
%%%% Define Theory%%%%
\newtheorem{definition}{\textbf{Definition}}
\newtheorem{theorem}{Theorem}
\newtheorem{exmp}{Example}[section]

%%%% SELF_DEFINE TERMS %%%%
\newcommand{\system }{\emph{VSJoin}\xspace} 
\newcommand{\sageflow }{\emph{sageflow}\xspace} 
\newcommand{\dsc }{\emph{DSC}\xspace}
\newcommand{\dscf }{\emph{Data Stream Clustering}\xspace}
\newcommand{\dsclf }{\emph{Data Stream Classification}\xspace}
\newcommand{\cef }{\emph{Concept Evolution}\xspace}

%%%% Define Design Aspects%%%%
\newcommand{\window }{\emph{Window Model}\xspace}
\newcommand{\outlier }{\emph{Outlier Detection}\xspace}
\newcommand{\data }{\emph{Data Structure}\xspace}
\newcommand{\offline }{\emph{Offline Refinement}\xspace}

%%%% Define Algorithm %%%% 
% \newcommand{\skm}{\hyperref[algo:streamkm++]{\underline{\textsf{StreamKM++}}}\xspace}
% \newcommand{\birch}{\hyperref[algo:birch]{\underline{\textsf{BIRCH}}}\xspace}
% \newcommand{\dstream}{\hyperref[algo:dstream]{\underline{\textsf{DStream}}}\xspace}
% \newcommand{\dbstream}{\hyperref[algo:dbstream]{\underline{\textsf{DBStream}}}\xspace}
% \newcommand{\denstream}{\hyperref[algo:denstream]{\underline{\textsf{DenStream}}}\xspace}
% \newcommand{\edmstream}{\hyperref[algo:edmstream]{\underline{\textsf{EDMStream}}}\xspace}
% \newcommand{\clustream}{\hyperref[algo:clustream]{\underline{\textsf{CluStream}}}\xspace}
% \newcommand{\slkmeans}{\hyperref[algo:slkmeans]{\underline{\textsf{SL-KMeans}}}\xspace}
% \newcommand{\kmeans}{\textit{{KMeans}}\xspace}
% \newcommand{\kmeanspp}{\textit{{KMeans++}}\xspace}
% \newcommand{\dbscan}{\textit{{DBSCAN}}\xspace}

\newcommand{\skm}{\textsf{StreamKM++}\xspace}
\newcommand{\birch}{\textsf{BIRCH}\xspace}
\newcommand{\dstream}{\textsf{DStream}\xspace}
\newcommand{\dbstream}{\textsf{DBStream}\xspace}
\newcommand{\denstream}{\textsf{DenStream}\xspace}
\newcommand{\edmstream}{\textsf{EDMStream}\xspace}
\newcommand{\clustream}{\textsf{CluStream}\xspace}
\newcommand{\slkmeans}{\textsf{SL-KMeans}\xspace}
\newcommand{\kmeans}{\textsf{KMeans}\xspace}
\newcommand{\kmeanspp}{\textsf{KMeans++}\xspace}
\newcommand{\dbscan}{\textsf{DBSCAN}\xspace}

%%%% Define DataSet %%%%
\newcommand{\fct}{\emph{FCT}\xspace}
\newcommand{\kdd}{\emph{KDD99}\xspace}
\newcommand{\sensor}{\emph{Sensor}\xspace}
\newcommand{\insect}{\emph{Insects}\xspace}
\newcommand{\eds}{\emph{EDS}\xspace}
\newcommand{\ods}{\emph{ODS}\xspace}

%%%% Define Structure %%%%
\newcommand{\CFT}{\hyperref[subsec:cf-tree]{\underline{\emph{CF Tree}}}\xspace}
\newcommand{\MCs}{\hyperref[subsec:MCs]{\underline{\emph{MCs}}}\xspace}
\newcommand{\CoreT}{\hyperref[subsec:CFT]{\underline{\emph{Coreset Tree}}}\xspace}
\newcommand{\Grids}{\hyperref[subsec:Grids]{\underline{\emph{Grids}}}\xspace}
\newcommand{\DPT}{\hyperref[subsec:DPT]{\underline{\emph{DP Tree}}}\xspace}
\newcommand{\MS}{\hyperref[subsec:meyerson_sketch]{\underline{\emph{Meyerson Sketch}}}\xspace}

%%%% Define Window %%%%
\newcommand{\landmark}{\hyperref[subsec:landmark]{\underline{\emph{Landmark Window Model}}}\xspace}
\newcommand{\sliding}{\hyperref[subsec:sliding]{\underline{\emph{Sliding Window Model}}}\xspace}
\newcommand{\damped}{\hyperref[subsec:damped]{\underline{\emph{Damped Window Model}}}\xspace}

%%%% Define Outliers %%%%
\newcommand{\buffer}{\hyperref[subsec:buffer]{\underline{\emph{OD w/ Buffer}}}\xspace}
\newcommand{\timer}{\hyperref[subsec:timer]{\underline{\emph{OD w/ Timer}}}\xspace}

%%%% Offline %%%%
\newcommand{\refinement}{\hyperref[sec:offline_refinement_strategy]{\underline{\emph{Offline Refinement}}}\xspace}

\usepackage{hyphenat}
\hyphenpenalty=10000
\tolerance=1000
\sloppy

\newcommand{\compact}{\vspace{-5pt}}
\newcommand{\subcompact}{\vspace{-6pt}}

\newcommand{\bib}{{\textcolor{red}{\tony{~\cite{xxx}}}}}

% \usepackage[skip=1pt]{caption}
% \setlength{\belowcaptionskip}{-5pt}

\renewcommand{\algorithmicrequire}{\textbf{Input:}} %Use Input in the format of Algorithm
\renewcommand{\algorithmicensure}{\textbf{Output:}} %UseOutput in the format of Algorithm

\newcounter{mycounter}
\setcounter{mycounter}{1}
\usepackage{comment}

%%%% Define Table %%%%

\newcommand{\notationTable}{{
\begin{table}[t]
\caption{Used Notations}
\begin{tabular}{|l|l|}
\hline
Notations   & Description                             \\ \hline
$DSC$       & Data Stream clustering  \\ \hline
$t$         & Timestamp of input data       \\ \hline
$w$         & Weight of input data       \\ \hline
$x$         & An input data object                    \\ \hline
$S$         & An input data stream                    \\ \hline
$d$         & Dimension of input data point                    \\ \hline
$c$         & Predicted cluster sets\\ \hline
% $GT$        & Ground truth cluster sets\\ \hline
% $F$         & False clustering point sets\\ \hline
% $h$         & Window length (msec)                    \\ \hline
$n$         & Length of current data stream                    \\ \hline
% $v$         & Input arrival rate (data objects/msec)  \\ \hline
$CF$        & Clustering Feature  \\ \hline
$MC$        & Micro cluster  \\ \hline

%$\lambda$             & Decay factor                            \\ \hline
\end{tabular}%
\label{tab:notation}
\end{table}
}}

\newcommand{\algorithmTable}{{
\begin{table*}[]
\centering
\caption{A summary of representative \dsc algorithms and their design decisions. \textmd{The year attribute for each algorithm is when it was first published.}}
\label{tab:summary}
\resizebox{0.8\textwidth}{!}{%
\begin{tabular}{|l|l|ll|l|l|l|}
\hline
\multirow{2}{*}{Algorithm}                                                                            & \multirow{2}{*}{Year} & \multicolumn{2}{c|}{Summarizing Data Structure}             & \multirow{2}{*}{Window Model} & \multirow{2}{*}{Outlier Detection Mechanism} & \multirow{2}{*}{Offline Clustering} \\ \cline{3-4}
                                                                                                      &                       & \multicolumn{1}{l|}{Name}                    & Catalog      &                               &                                    &                                     \\ \hline
{\color{blue} \birch}~\cite{BIRCH:96}        & 1996                  & \multicolumn{1}{l|}{Clustering Feature Tree} & Hierarchical & Landmark                      & w/o buffer, w/o timer                   & N.A.                                \\ \hline
{\color{blue} \clustream}~\cite{Clustream:03}& 2003                  & \multicolumn{1}{l|}{Micro Clusters}          & Partition    & Landmark                      & w/o buffer, w/ timer                       & Distance-based                      \\ \hline
{\color{blue} \denstream}~\cite{cao2006density}& 2006                  & \multicolumn{1}{l|}{Micro Clusters}          & Partition    & Damped                        & w/ buffer, w/ timer                      & Density-based                       \\ \hline
{\color{blue} \dstream}~\cite{DStream:2007}  & 2007                  & \multicolumn{1}{l|}{Grid}                    & Partition    & Damped                        & w/o buffer, w/ timer                    & N.A.                                \\ \hline
{\color{blue} \skm}~\cite{StreamKM++:12}     & 2012                  & \multicolumn{1}{l|}{Coreset Tree}            & Hierarchical & Landmark                      & N.A.                               & Distance-based                      \\ \hline
{\color{blue} \dbstream}~\cite{DBStream:16}  & 2016                  & \multicolumn{1}{l|}{Micro Clusters}          & Partition    & Damped                        & w/o buffer, w/ timer                       & Density-based                       \\ \hline
{\color{blue} \edmstream}~\cite{EDMStream:17}& 2017                  & \multicolumn{1}{l|}{Dependency Tree}            & Hierarchical & Damped                        & w/ buffer, w/ timer                   & N.A.                               \\ \hline
{\color{blue} \slkmeans}~\cite{SL-KMeans:20} & 2020                  & \multicolumn{1}{l|}{Meyerson Sketch}         & Hierarchical & Sliding                       & N.A.                               & N.A.                                \\ \hline
\end{tabular}%
% \begin{tabular}{|l|l|ll|l|l|l|}
% \hline
% \multirow{2}{*}{\textbf{Algorithm}}                                                                            & \multirow{2}{*}{\textbf{Year}} & \multicolumn{2}{c|}{\textbf{Summarizing Data Structure}}                       & \multirow{2}{*}{\textbf{Window Model}} & \multirow{2}{*}{\textbf{Outlier Detection}} & \multirow{2}{*}{\textbf{Offline Refinement}} \\ \cline{3-4}
%                                                                                                       &                       & \multicolumn{1}{l|}{\emph{Primitive Data Structure}} & \emph{Global Data Structure} &                               &                                    &                                     \\ \hline
% {\color{blue} \birch}~\cite{BIRCH:96}         & 1996                  & \multicolumn{1}{l|}{Clustering feature}       & Tree          & Landmark                      & w/o buffer, w/o timer                     & N.A.                                \\ \hline
% {\color{blue} \clustream}~\cite{Clustream:03} & 2003                  & \multicolumn{1}{l|}{Micro cluster}            & List             & Landmark                      & w/o buffer, w/ timer                  & Distance-based                      \\ \hline
% {\color{blue} \denstream}~\cite{DenStream:06} & 2006                  & \multicolumn{1}{l|}{Micro cluster}            & List             & Damped                        & w/ buffer, w/ timer                         & Density-based                       \\ \hline
% {\color{blue} \dstream}~\cite{DStream:2007}   & 2007                  & \multicolumn{1}{l|}{Grid}                     & List             & Damped                        & w/o buffer, w/ timer                       & N.A.                                \\ \hline
% {\color{blue} \skm}~\cite{StreamKM++:12}      & 2012                  & \multicolumn{1}{l|}{Raw point}                & Tree          & Landmark                      & N.A.                               & Distance-based                      \\ \hline
% {\color{blue} \dbstream}~\cite{DBStream:16}   & 2016                  & \multicolumn{1}{l|}{Micro cluster}            & List             & Damped                        & w/o buffer, w/ timer                     & Density-based                       \\ \hline
% {\color{blue} \edmstream}~\cite{EDMStream:17} & 2017                  & \multicolumn{1}{l|}{Cluster Cell}             & Tree          & Damped                        & w/ buffer, w/ timer                         & N.A.                                \\ \hline
% {\color{blue} \slkmeans}~\cite{SL-KMeans:20}  & 2020                  & \multicolumn{1}{l|}{Raw point}                & List             & Sliding                       & N.A.                               & N.A.                                \\ \hline
% \end{tabular}%
}
\end{table*}
}}

\newcommand{\workloadTable}{{
\begin{table}[t]
\centering
\caption{Characteristics differences of selected workloads. \textmd{Note that the outliers column refers to whether there are outliers in the final clustering results.}}
\label{tab:Workload}
\resizebox{0.47\textwidth}{!}{%
\begin{tabular}{|c|c|c|c|c|}
\hline
Workload                                    & Length                & Dimension                 & Outliers             & Evolving Frequency\\ \hline
{\color{blue}\fct}~\cite{Covertype}         & 581012                & 54                      & False                          & Low\\ \hline
{\color{blue}\kdd}~\cite{KDD99}             & 494021                & 41                      & True                         & Low\\ \hline
{\color{blue}\insect}~\cite{Insect}         & 905145                & 33                      & False                         & Low\\ \hline
{\color{blue}\sensor}~\cite{SensorData}     & 2219803               & 5                       & False                         & High\\ \hline
{\color{blue}\eds}~\cite{DenStream:06}      & 245270                & 2                       & False                         & Varying\\\hline
{\color{blue}\ods}                          & 100000                & 2                       & Varying                       & High\\ \hline
\end{tabular}%
}
\end{table}
}}


\newcommand{\designTable}{{
% Please add the following required packages to your document preamble:
\begin{table*}[t]
\centering
\caption{Summary of Design Aspect Assembly Versions based on Different Design Choices}
\label{tab:design}
\resizebox{0.99\textwidth}{!}{%
\begin{tabular}{|l|l|l|ll|l|}
\hline
\multirow{2}{*}{Version} & \multirow{2}{*}{Window Model} & \multirow{2}{*}{Outlier Detection} & \multicolumn{2}{c|}{Data Structure}                                                                                 & \multirow{2}{*}{Offline Clustering} \\ \cline{4-5}
                           &                               &                                    & \multicolumn{1}{l|}{Primitive Data Structure} & \multicolumn{1}{l|}{Global Data Structure}   &                                     \\ \hline
{\color{blue}V1}           & Landmark          & Buffer-based            & \multicolumn{1}{l|}{Micro cluster}         & \multicolumn{1}{l|}{Tree}          & Distance-based                    \\ \hline
{\color{blue}V2}           & Landmark          & Buffer-based          & \multicolumn{1}{l|}{Micro cluster}         & \multicolumn{1}{l|}{Tree}          & Density-based                    \\ \hline
{\color{blue}V3}           & Landmark          & Buffer-based           & \multicolumn{1}{l|}{Micro cluster}         & \multicolumn{1}{l|}{Tree}         & N.A                       \\ \hline
{\color{blue}V4}           & Sliding           & Buffer-based            & \multicolumn{1}{l|}{Micro cluster}         & \multicolumn{1}{l|}{Tree}          & N.A                        \\ \hline
{\color{blue}V5}           & Damped            & Buffer-based           & \multicolumn{1}{l|}{Micro cluster}         & \multicolumn{1}{l|}{Tree}          & N.A                    \\ \hline
{\color{blue}V6}           & Landmark          & Distance-based                      & \multicolumn{1}{l|}{Micro Cluster}         & \multicolumn{1}{l|}{Tree}         & N.A.                                  \\ \hline
{\color{blue}V7}           & Landmark          & Density-based             & \multicolumn{1}{l|}{Micro cluster}         & \multicolumn{1}{l|}{Tree}          & N.A                    \\ \hline
{\color{blue}V8}           & Landmark          & Buffer-based            & \multicolumn{1}{l|}{Micro cluster}         & \multicolumn{1}{l|}{List}             & N.A.                  \\ \hline
{\color{blue}V9}           & Landmark          & Buffer-based            & \multicolumn{1}{l|}{Grid}                  & \multicolumn{1}{l|}{List}             & N.A.                   \\ \hline
\end{tabular}%
}
\end{table*}
}}